% Options for packages loaded elsewhere
\PassOptionsToPackage{unicode}{hyperref}
\PassOptionsToPackage{hyphens}{url}
%
\documentclass[
  11pt,
  letterpaper,
]{article}
\usepackage{amsmath,amssymb}
\usepackage{iftex}
\ifPDFTeX
  \usepackage[T1]{fontenc}
  \usepackage[utf8]{inputenc}
  \usepackage{textcomp} % provide euro and other symbols
\else % if luatex or xetex
  \usepackage{unicode-math} % this also loads fontspec
  \defaultfontfeatures{Scale=MatchLowercase}
  \defaultfontfeatures[\rmfamily]{Ligatures=TeX,Scale=1}
\fi
\usepackage{lmodern}
\ifPDFTeX\else
  % xetex/luatex font selection
\fi
% Use upquote if available, for straight quotes in verbatim environments
\IfFileExists{upquote.sty}{\usepackage{upquote}}{}
\IfFileExists{microtype.sty}{% use microtype if available
  \usepackage[]{microtype}
  \UseMicrotypeSet[protrusion]{basicmath} % disable protrusion for tt fonts
}{}
\makeatletter
\@ifundefined{KOMAClassName}{% if non-KOMA class
  \IfFileExists{parskip.sty}{%
    \usepackage{parskip}
  }{% else
    \setlength{\parindent}{0pt}
    \setlength{\parskip}{6pt plus 2pt minus 1pt}}
}{% if KOMA class
  \KOMAoptions{parskip=half}}
\makeatother
\usepackage{xcolor}
\usepackage[margin=1in]{geometry}
\usepackage{longtable,booktabs,array}
\usepackage{calc} % for calculating minipage widths
% Correct order of tables after \paragraph or \subparagraph
\usepackage{etoolbox}
\makeatletter
\patchcmd\longtable{\par}{\if@noskipsec\mbox{}\fi\par}{}{}
\makeatother
% Allow footnotes in longtable head/foot
\IfFileExists{footnotehyper.sty}{\usepackage{footnotehyper}}{\usepackage{footnote}}
\makesavenoteenv{longtable}
\setlength{\emergencystretch}{3em} % prevent overfull lines
\providecommand{\tightlist}{%
  \setlength{\itemsep}{0pt}\setlength{\parskip}{0pt}}
\setcounter{secnumdepth}{-\maxdimen} % remove section numbering
\newlength{\cslhangindent}
\setlength{\cslhangindent}{1.5em}
\newlength{\csllabelwidth}
\setlength{\csllabelwidth}{3em}
\newlength{\cslentryspacingunit} % times entry-spacing
\setlength{\cslentryspacingunit}{\parskip}
\newenvironment{CSLReferences}[2] % #1 hanging-ident, #2 entry spacing
 {% don't indent paragraphs
  \setlength{\parindent}{0pt}
  % turn on hanging indent if param 1 is 1
  \ifodd #1
  \let\oldpar\par
  \def\par{\hangindent=\cslhangindent\oldpar}
  \fi
  % set entry spacing
  \setlength{\parskip}{#2\cslentryspacingunit}
 }%
 {}
\usepackage{calc}
\newcommand{\CSLBlock}[1]{#1\hfill\break}
\newcommand{\CSLLeftMargin}[1]{\parbox[t]{\csllabelwidth}{#1}}
\newcommand{\CSLRightInline}[1]{\parbox[t]{\linewidth - \csllabelwidth}{#1}\break}
\newcommand{\CSLIndent}[1]{\hspace{\cslhangindent}#1}
\ifLuaTeX
\usepackage[bidi=basic]{babel}
\else
\usepackage[bidi=default]{babel}
\fi
\babelprovide[main,import]{american}
% get rid of language-specific shorthands (see #6817):
\let\LanguageShortHands\languageshorthands
\def\languageshorthands#1{}
\ifLuaTeX
  \usepackage{selnolig}  % disable illegal ligatures
\fi
\IfFileExists{bookmark.sty}{\usepackage{bookmark}}{\usepackage{hyperref}}
\IfFileExists{xurl.sty}{\usepackage{xurl}}{} % add URL line breaks if available
\urlstyle{same}
\hypersetup{
  pdftitle={Canonical witnesses for modular convex subsequences with a small index deficit},
  pdfauthor={Stamatis-Christos Saridakis},
  pdflang={en-US},
  hidelinks,
  pdfcreator={LaTeX via pandoc}}

\title{Canonical witnesses for modular convex subsequences with a small
index deficit}
\author{Stamatis-Christos Saridakis}
\date{24 October 2026}

\begin{document}
\maketitle

\textbf{Research draft.} This technical note states a specialized upper
bound for a corrected research implementation. Literature priority and
suitability for a publication venue remain unresolved.

\hypertarget{abstract}{%
\section{Abstract}\label{abstract}}

Consider a maximum-cardinality subsequence of an integer array whose
original indices sum to zero modulo a prescribed integer and whose
consecutive selected values have strictly increasing differences. We
study the regime where the modulus is close to the sum of all indices.
Writing \(h=n(n+1)/2-m<n\) and \(H=h+1\), we prove that the globally
lexicographically least entire optimal index tuple can be computed in
\(O(n+H^{5/2}\log^2(h+2))\) wide-word operations using \(O(H^2)\)
auxiliary words, excluding the input and one returned output. The
algorithm combines established sorted-difference dynamic programming and
canonical self-reduction with a deficit-specific reduction. The omitted
indices sum to \(h\), giving a core of at most \(h+2\) indices,
\(O(H\log(h+2))\) eligible transitions, and at most \(O(\sqrt H)\)
omissions. A forced-prefix scalar oracle and monotone threshold queries
reconstruct the complete canonical witness without storing witness
histories. An accompanying computational audit checks 80,283
complete-witness instances and identifies a linear auxiliary allocation
in the original implementation's final validation; an indexed validation
and deterministic output cursor repair it. The contribution is a
specialized composition with explicit complexity accounting, without a
claim of a new general dynamic-programming primitive or established
literature priority.

\hypertarget{introduction}{%
\section{1 Introduction}\label{introduction}}

An optimal score and a canonical optimal witness impose different
obligations on an exact algorithm. For a convex subsequence, a scalar
score can be propagated in sorted order of value differences. Retaining
one locally preferred predecessor does not by itself show that the
resulting entire index tuple is globally lexicographically least. A
second difficulty appears in space analysis: even when the optimization
state is small, copying a large returned witness during validation can
invalidate an auxiliary-space bound.

This note addresses both obligations for a modular index constraint in a
restricted parameter regime. The input is an integer array \(A[1..n]\).
Selected values must be strictly convex in their selected order, and the
sum of the selected original indices must be divisible by \(m\). The
objective first maximizes the number of selected indices and then
minimizes their complete increasing tuple lexicographically. If \(m\) is
sufficiently large, the modular condition becomes an exact
omitted-weight condition. When its deficit \(h\) is smaller than \(n\),
all choices occur near the start of the array.

The main result is a complete algorithm and proof with quadratic
auxiliary space in \(h+1\). Its central reconstruction step selects
omissions rather than testing every input position: an optimal omitted
set has only \(O(\sqrt{h+1})\) elements, and each next omission is found
by binary search on a threshold feasibility predicate. The proof
accounts for every repeated query, the compulsory suffix, event sorting,
and output construction.

The mathematical mechanisms have clear antecedents. Accordingly, the
intended contribution is a specialized bound and a reproducible audit of
its implementation. The manuscript does not claim a general advance in
canonical traceback, a lower-bound-optimal algorithm, or a demonstrated
practical speedup. The corrected artifact and the original
implementation are preserved separately in the MOCOSO research
repository (Saridakis 2026).

\hypertarget{problem-and-computational-model}{%
\section{2 Problem and computational
model}\label{problem-and-computational-model}}

Let \(n\geq2\), let \(A[1..n]\) contain arbitrary signed integers, and
let \(m\geq2\) be an integer. Define

\[U=\frac{n(n+1)}2,\qquad h=U-m,\qquad H=h+1.\]

Assume \(U/2<m\leq U\) and \(0\leq h<n\). A feasible witness is a
nonempty tuple \(S=(s_1,\ldots,s_k)\) with strictly increasing entries
in \([n]=\{1,\ldots,n\}\) satisfying

\[\sum_{r=1}^k s_r\equiv0\pmod m,\]

and, for each \(1\leq r\leq k-2\),

\[A[s_{r+2}]-A[s_{r+1}]>A[s_{r+1}]-A[s_r].\]

Tuples of length one or two satisfy the convexity condition vacuously.
This is convexity of selected values in selected order. It is not a
condition on divided differences using the original index gaps. Among
all feasible witnesses, the requested output has maximum length and,
among witnesses of that length, is lexicographically least. An empty
tuple denotes infeasibility and is not a feasible candidate of score
zero.

We count operations on wide words large enough to contain input values,
differences, triangular sums, addresses, and polynomially bounded
counters. If signed input values require at most \(B\) bits,
\(b=\Theta(B+\log(n+2))\) bits per word suffice, with an adequate
constant factor. Arithmetic and comparisons on such words have unit cost
in the primary theorem. This is not a fixed 64-bit RAM assumption. The
auxiliary-space measure excludes a read-only random-access input and one
final output array; all intermediate arrays, event records, sorting
storage, and temporary output copies are charged.

\textbf{Theorem 1.} Under these assumptions, the globally
lexicographically least entire maximum-cardinality feasible witness, or
an infeasibility report, is computable deterministically in
\(O(n+H^{5/2}\log^2(h+2))\) wide-word operations and \(O(H^2)\)
auxiliary wide words.

The theorem concerns the algorithm below and the corrected research
derivative described in Section 7. The original pinned implementation
has a separate auxiliary-space defect. With elementary bit arithmetic, a
conservative translation of the theorem gives time
\(O((n+H^{5/2}\log^2(h+2))b^2)\) and auxiliary storage \(O(H^2b)\) bits.
The excluded input and output require \(O(nB)\) and \(O(n\log(n+2))\)
bits, respectively. These bounds do not predict Python process RSS or
allocator behavior.

\hypertarget{established-ingredients-and-scope-of-contribution}{%
\section{3 Established ingredients and scope of
contribution}\label{established-ingredients-and-scope-of-contribution}}

The unweighted value-convex subsequence problem appears in Erickson's
algorithms textbook (Erickson 2019, 126, exercise (f)). Kluev's 2013
public technical answer gives the closely related sorted-difference
sweep, including decreasing destination indices for equal differences,
scalar endpoint scores, and witness reconstruction suggestions (Kluev
2013, steps 1--6). That answer is a primary public disclosure, not a
peer-reviewed paper. The sweep and its equal-difference ordering are
therefore established ingredients, not claimed contributions of this
note.

Exact resource coordinates in dynamic programming are also standard.
Turner describes equality-constrained resource shortest paths on acyclic
digraphs and gives a pseudopolynomial algorithm (Turner 2011, equation
(23), Theorem 2). Her result supplies an appropriate resource-DP
template, but not the compressed event invariant or the
deficit-dependent transition count needed here. Simply expanding a
last-two-indices graph would not establish our stated bound.

Canonical witness search by constrained extension queries belongs to
established self-reduction methods; Große, Rothe, and Wechsung
explicitly discuss lexicographically first solutions in that setting
(Große, Rothe, and Wechsung 2006, sec. 3). Our use of threshold queries
needs a proof that the particular constrained optimization oracle
remains efficient. Unconstrained decision alone is insufficient for that
conclusion. Low-space exact resource algorithms also occur in Subset Sum
(Jin, Vyas, and Williams 2020); their problem does not include
selected-order convexity or this entire-index-tuple objective, so no
black-box implication is asserted here.

The surviving contribution is the joint specialized analysis: an exact
omission deficit, a two-anchor core, a harmonic eligible-event bound, a
constrained scalar sweep, and sparse omission reconstruction with all
queries charged. These statements can be derived by combining familiar
mechanisms with elementary problem structure. We did not identify the
same complete theorem in the sources inspected on 24 October 2026. That
bounded observation does not establish firstness, and an earlier
equivalent result would substantially weaken the paper's originality
claim.

\hypertarget{deficit-structure-and-eligible-transitions}{%
\section{4 Deficit structure and eligible
transitions}\label{deficit-structure-and-eligible-transitions}}

\textbf{Lemma 2 (exact complement).} For any feasible witness, its
omitted set \(D\), consisting of the indices outside \(S\), satisfies
\(\sum_{d\in D}d=h\). Conversely, any complement with that omitted sum
satisfies the modular index condition. Every omission is at most \(h\),
and the omission count \(t\) satisfies \(t(t+1)/2\leq h\).

\emph{Proof.} A nonempty selected sum is positive and at most \(U<2m\).
Its only possible positive multiple of \(m\) is \(m\). Taking
complements gives omitted sum \(U-m=h\), and the converse follows
immediately. An omitted index greater than \(h\) would already exceed
the nonnegative total. Distinct positive omitted indices have sum at
least \(1+\cdots+t=t(t+1)/2\). Thus \(t=O(\sqrt H)\). In particular,
every index from \(h+1\) through \(n\) is compulsory. \(\square\)

Set \(p=\min(n,h+2)\) and let \(W(q)=q(q+1)/2\) for \(q\geq0\).

\textbf{Lemma 3 (two-anchor core).} Test the original consecutive
triples ending at indices \(k=h+3,\ldots,n\). If any fails strict
convexity, the instance is infeasible. Otherwise, feasible core
witnesses on \([p]\) that include its compulsory indices correspond
exactly to full feasible witnesses by appending \(p+1,\ldots,n\). Scores
increase by \(n-p\).

\emph{Proof.} Each tested triple consists entirely of compulsory indices
and hence occurs consecutively in every witness. When \(p<n\), the core
contains the two compulsory indices \(h+1,h+2\). Consequently, the first
triple crossing into the appended suffix ends at \(h+3\) and is among
those tested. Every later suffix triple is tested as well. Any triple
involving a mutable index and ending at either compulsory anchor remains
in the core. If \(p=n\), there is no suffix to append. Restriction and
extension therefore preserve feasibility in both directions. \(\square\)

The suffix check costs \(O(n)\) once. Including only \(h+1\) as an
anchor would leave a triple involving the last mutable selected index
unchecked; the second anchor is essential to this reduction.

A core event is a possible pair of consecutive selected indices
\((i,j)\) with \(1\leq i<j\leq p\). Retain the pair only if its open gap
contains no compulsory index and its omitted weight

\[g(i,j)=W(j-1)-W(i)\]

is at most \(h\). Adjacent indices have gap weight zero. Every
consecutive selected pair of a feasible core witness is eligible. For
each retained event, store its difference \(d(i,j)=A[j]-A[i]\) and its
gap weight.

\textbf{Lemma 4 (harmonic event count).} The number \(M\) of eligible
events is \(O(H\log(h+2))\). They can be enumerated in \(O(M+p)\) time
and sorted in \(O(M\log(M+1))\) time.

\emph{Proof.} There are at most \(p-1\) adjacent events. For a nonempty
gap containing \(u=j-i-1\) indices, its weight is

\[g(i,j)=ui+\frac{u(u+1)}2.\]

Eligibility implies \(u\leq\operatorname{floor}(\sqrt{2h})\) and
\(i\leq\operatorname{floor}(h/u)\). Hence

\[M\leq p-1+\sum_{u=1}^{\operatorname{floor}(\sqrt{2h})}\operatorname{floor}(h/u)
=O(H\log(h+2)).\]

This is an upper bound, not an assertion that every counted pair is
eligible. For each fixed \(i\), increasing \(j\) increases the gap
weight and cannot remove a compulsory-index violation. Thus a direct
scan stops at its first failure, charging eligible events and at most
one failed candidate per \(i\). Comparison sorting gives the claimed
preprocessing cost. When \(h=0\), the sum is empty and the bounds remain
valid. \(\square\)

\hypertarget{a-forced-prefix-scalar-optimization-oracle}{%
\section{5 A forced-prefix scalar optimization
oracle}\label{a-forced-prefix-scalar-optimization-oracle}}

Fix a set \(Q\subseteq[h]\) of forced omissions and a threshold
\(x\in[1,h+1]\). The oracle \(F(Q,x)\) returns the maximum full-witness
length subject to the original constraints, omission of every member of
\(Q\), and selection of every index less than \(x\) outside \(Q\). It
returns zero if no such witness exists. The suffix check from Lemma 3 is
performed before any oracle calls.

Let a Boolean array mark forbidden endpoints in \(Q\), and define a
prefix count

\[R(v)=|\{r\leq v:r<x,\ r\notin Q\}|,\qquad R(0)=0.\]

An event \((i,j)\) is allowed for this query precisely when neither
endpoint is forbidden and \(R(j-1)=R(i)\), so its gap contains no
required index. These checks take constant time per event. A singleton
start at \(i\) is allowed when \(i\leq\min(p,h+1)\), \(i\notin Q\),
\(W(i-1)\leq h\), and \(R(i-1)=0\).

Create a matrix \(V[i,r]\) for \(1\leq i\leq p\) and \(0\leq r\leq h\),
initially unreachable. For each allowed singleton start, set
\(V[i,W(i-1)]=1\). Sort all eligible events once, outside the oracle, by
increasing difference and then decreasing destination. Remaining ties
may be broken arbitrarily. For each allowed event \((i,j)\), and for
each \(0\leq r\leq h-g(i,j)\) with reachable \(V[i,r]\), perform

\[V[j,r+g(i,j)]\gets\max\{V[j,r+g(i,j)],V[i,r]+1\}.\]

The result is \(V[p,h]+n-p\) if \(V[p,h]\) is reachable, and zero
otherwise. The matrix stores only lengths. No predecessor tuple or
persistent witness tree is retained.

\textbf{Lemma 5 (event-relative dominance).} Immediately before an event
\((i,j)\) of difference \(d\) is processed, the source values \(V[i,r]\)
represent optimal lengths among admissible paths ending at \(i\), of
omitted weight \(r\) through \(i\), whose last difference is less than
\(d\), together with admissible singleton starts. The update therefore
enforces strict convexity, and the final oracle value is correct.

\emph{Proof.} All events with difference less than \(d\) have been
processed. Any event of difference exactly \(d\) entering \(i\) has
destination \(i<j\), and decreasing destination order places it after
the current event. Thus no equal-difference incoming event can yet have
contributed to \(V[i,r]\). Events with larger difference also have not
been processed. Singleton starts have no last difference and may precede
any first edge.

Inductively, every reachable state arises from a strictly convex
selected path satisfying the endpoint and required-gap checks.
Conversely, the previous edge of any admissible path that can use
\((i,j)\) has smaller difference and has already been processed, so that
path's predecessor score is represented or dominated. At this event, a
larger score with the same endpoint and exact resource can use every
continuation available to the smaller score: the source paths all
satisfy the current difference threshold. Retaining only the maximum
score therefore loses no optimal extension. This reasoning is relative
to the event order; it does not claim that arbitrary endpoint-equal
paths with different last differences are interchangeable at all times.

Singleton omission weights and the disjoint open gaps count exactly the
indices omitted through the current endpoint. Forbidden endpoints ensure
that members of \(Q\) are absent, and required-prefix checks ensure that
every required index is present. Eligibility prevents compulsory indices
from being skipped, including by an initial start. Since \(p\) is
compulsory, an admissible core path must end at \(p\) and have total
omitted weight \(h\). Lemma 3 then supplies exactly the appended suffix
length. This proves both soundness and completeness. \(\square\)

\textbf{Lemma 6 (oracle resources).} A query costs \(O(H^2\log(h+2))\)
operations and \(O(H^2)\) auxiliary words.

\emph{Proof.} The matrix has \(O(pH)=O(H^2)\) entries. Its
initialization, flags, and prefix counts take \(O(H^2)\) time. Each of
\(M\) events scans at most \(H\) resource entries, for \(O(MH)\) time.
The event list occupies \(O(M)\) words and is reused. The matrix is
released between calls. Lemma 4 gives the bounds. There is no suffix or
output scan inside a query. \(\square\)

\hypertarget{canonical-reconstruction-and-total-complexity}{%
\section{6 Canonical reconstruction and total
complexity}\label{canonical-reconstruction-and-total-complexity}}

Compute \(L=F(\varnothing,1)\). If \(L=0\), report infeasibility.
Otherwise all optimal witnesses omit exactly \(t=n-L\) indices. Maintain
an increasing omission prefix \(Q=(q_1,\ldots,q_r)\) that is the
beginning of the omitted tuple of at least one optimum. Initially \(Q\)
is empty and its previous index \(b\) is zero. Until \(r=t\),
binary-search thresholds \(x\in[b+1,h+1]\) for the greatest value
satisfying

\[F(Q,x)=L.\]

Append that value to \(Q\) and update \(b\). Return the increasing
complement of the completed \(Q\).

\textbf{Lemma 7 (threshold monotonicity and extension).} With the stated
prefix invariant and \(r<t\), the predicate \(F(Q,x)=L\) is true exactly
when some optimal omitted tuple extending \(Q\) has its next element at
least \(x\). It is true at \(b+1\), false at \(h+1\), and monotone from
true to false. Its greatest true threshold is the greatest possible next
omission among those optima.

\emph{Proof.} Every index below \(x\) outside \(Q\) must be selected.
Since \(x>b\) and all members of \(Q\) are forbidden, this fixes the
omission prefix through \(x-1\) to exactly \(Q\). There remains at least
one omission because \(r<t\). It must lie at or beyond \(x\).
Conversely, an optimum with that omission prefix and next omission at
least \(x\) satisfies all query constraints. Constrained scores cannot
exceed the unconstrained optimum \(L\), so equality is exactly the
claimed existence test.

Increasing \(x\) adds required selections and cannot create an optimum.
The prefix invariant gives truth at \(b+1\). At \(h+1\), every mutable
index outside \(Q\) would be required, leaving no location for the
remaining omission. Thus the upper endpoint is false. At a greatest true
threshold \(x^*\), a witnessing optimum cannot have its next omission
greater than \(x^*\), because it would then satisfy threshold \(x^*+1\)
as well. Its next omission is therefore exactly \(x^*\). Appending it
preserves the prefix invariant and makes the next omission as large as
possible. \(\square\)

This predicate concerns a lower bound on the next omission, not equality
to a proposed omission. Equality feasibility need not be monotone. For
example, with \(A=(1,3,9,27,81,243,729)\), \(n=7\), and \(m=22\), the
maximum length is six and the only optimal omission is index six. The
equality predicate is false at one through five, true at six, and false
at seven. The threshold predicate is true at one through six and false
at seven.

\textbf{Lemma 8 (entire-witness order).} Lexicographically maximizing
the increasing omitted tuple among optimal witnesses lexicographically
minimizes the entire selected tuple.

\emph{Proof.} Compare distinct omitted sets of equal cardinality, and
let \(z\) be the smallest index in their symmetric difference. Their
increasing omitted tuples first differ in favor of the set that contains
\(z\): that tuple is lexicographically smaller. The selected complements
first differ in the reverse direction, since exactly the other
complement contains \(z\). Thus complementation reverses the order.
Lemma 7 successively maximizes every omission coordinate among surviving
optima, so the completed omitted tuple is globally greatest. Its
complement is globally least. \(\square\)

\emph{Proof of Theorem 1.} Correctness follows from Lemmas 2--8. For an
explicit query bound, count the initial optimization, one lower-endpoint
assertion query per omission, and at most
\(\operatorname{ceil}(\log_2 H)\) binary-search queries per omission.
This yields

\[q\leq1+t\bigl(1+\operatorname{ceil}(\log_2 H)\bigr).\]

For \(h=0\), \(t=0\), so only the initial query is needed. In all cases,
Lemma 2 gives \(q=O(1+\sqrt H\log(h+2))\). Multiplying by Lemma 6 gives
the claimed parameter-dependent time. Sorting the single event list
costs \(O(H\log^2(h+2))\), which is absorbed. The initial input and
suffix checks and final output construction cost \(O(n)\) in total; they
are not repeated by the oracle.

At any instant there is one \(O(H^2)\) matrix, an \(O(H\log(h+2))\)
event list with sorting workspace of the same order, \(O(H)\) flags and
prefix counters, and \(O(\sqrt H)\) chosen omissions. These occupy
\(O(H^2)\) auxiliary words. Allocate the known-length output once and
fill it by scanning \([n]\) against a monotone cursor in the sorted
omission list. Validate it by indexed reads with a constant number of
temporary words. No additional output-sized tuple or list is necessary.
This establishes the stated auxiliary-space bound. \(\square\)

\hypertarget{implementation-defect-and-correction}{%
\section{7 Implementation defect and
correction}\label{implementation-defect-and-correction}}

The original G9 source is pinned to repository commit \texttt{96ac7e2}
(full identifiers are in the evidence manifest). Its final convexity
assertion iterates over
\texttt{zip(selected,\ selected{[}1:{]},\ selected{[}2:{]})}. Python
evaluates both slices before the iteration, creating two additional
tuples of length proportional to the output. These tuples are auxiliary
storage even though the original returned tuple is excluded.

The counterfamily is \(h=0\), \(m=U\), and \(A[i]=i^2\). Every index
must be selected and the complete array is strictly convex. Therefore
\(H=1\) while the two slices use \(\Theta(n)\) auxiliary reference
slots. The original code's asserted \(O(H^2)\) auxiliary-space bound is
false on this family. Passing functional tests cannot detect that
asymptotic defect without a space-specific check.

The corrected derivative replaces the sliced assertion with indexed
comparisons and constructs the output through an omission cursor. The
cursor also removes reliance on average-case hash membership for a
deterministic output-time claim. Its iterator exposes the known
remaining output length to CPython; interpreter allocation details
remain outside the abstract storage theorem. The scalar DP, event order,
and reconstruction loop are unchanged. The derivative is an attributed
correction, not a clean-room reimplementation of G9.

Both versions and all prior research implementations are retained in
snapshot \texttt{43353e2} (Saridakis 2026). The corrected file is
\texttt{research/v03/g9\_independent\_audit/solver\_corrected.py}. Its
complete SHA-256 and the audit harness hash appear in the companion
evidence manifest, avoiding ambiguity about which implementation
supports the theorem.

\hypertarget{computational-audit}{%
\section{8 Computational audit}\label{computational-audit}}

The computational evidence supplements the proof. The preserved audit
uses seed 2026102415 and CPython 3.12.14 on Linux x86-64. Three
separately written reference approaches were used: direct
selected-subset enumeration, full-array dynamic programming indexed by
the last two selected indices and exact selected sum, and distinct-part
omission enumeration for planted instances. The reference methods do not
import upstream reference solvers. Independence here means separately
written methods within the same assisted audit session, not blinded
external replication or formal verification.

The full-witness corpus contains 80,283 instances. The following
categories are disjoint; the additional pair-DP checks are a subset of
the random category.

\begin{longtable}[]{@{}
  >{\raggedright\arraybackslash}p{(\columnwidth - 4\tabcolsep) * \real{0.3333}}
  >{\raggedleft\arraybackslash}p{(\columnwidth - 4\tabcolsep) * \real{0.3333}}
  >{\raggedright\arraybackslash}p{(\columnwidth - 4\tabcolsep) * \real{0.3333}}@{}}
\toprule\noalign{}
\begin{minipage}[b]{\linewidth}\raggedright
Full-witness category
\end{minipage} & \begin{minipage}[b]{\linewidth}\raggedleft
Instances
\end{minipage} & \begin{minipage}[b]{\linewidth}\raggedright
Scope
\end{minipage} \\
\midrule\noalign{}
\endhead
\bottomrule\noalign{}
\endlastfoot
Exhaustive selected-subset comparisons & 73,809 & Three values, lengths
2 through 8, valid deficits \\
Seeded random comparisons & 6,000 & Lengths 3 through 12, six input
families \\
Planted omission instances & 450 & Deficits 0 through 34 \\
Maximum-omission stress cases & 6 & Up to 16 omissions \\
Constructed deep lexicographic ties & 18 & Shared omission prefixes of
lengths 0 through 5 \\
Total & 80,283 & Complete returned tuples compared \\
\end{longtable}

The exhaustive alphabet is \(\{-3,0,5\}\). Random families include small
signed values, perturbed quadratics, random walks, and very large
integers, including 120-digit magnitudes. Of the 6,000 random instances,
1,800 also receive the independent last-two-indices DP check. Across the
full corpus, 75,621 instances are infeasible and 4,662 feasible. This
imbalance limits any interpretation as a representative workload; the
suite is designed for correctness falsification, not performance
estimation.

The constrained oracle is tested separately on 13,767
forced-omission/threshold queries against direct enumeration; 1,565 of
those also receive the independent pair-DP check. Event enumeration and
its harmonic bound are checked for all deficits from zero through 200.
These counts are additional checks, not additional full-witness
instances. No mismatch is recorded for the original or corrected exact
returned witnesses in the preserved run.

Mutation tests deliberately reverse the equal-difference destination
order, remove forbidden-endpoint or required-prefix checks, and shorten
the core to one anchor. Concrete counterexamples in the audit artifact
show that these changes break the relevant invariants. For instance,
with all four values zero and \(m=8\), ascending equal-difference order
can return \((1,3,4)\) even though no feasible witness exists. The
mutant's assertions are disabled only to observe that invalid result;
assertions remain enabled for the normal audit.

An additional search over 10,000 random inputs found no hole in the
feasible interval of first omissions under the investigated hypothesis.
That failed falsification attempt establishes no interval theorem and is
not used in the proof. The proof uses the nested threshold constraints
directly. All explicit mutation examples and the failed-search scope are
preserved in the audit report.

The validation-phase allocation probe isolates the discovered slice
defect. After the returned tuple is constructed, it resets the tracing
peak and measures extra allocations during validation and return. It
does not measure total solver memory or RSS.

\begin{longtable}[]{@{}rrr@{}}
\toprule\noalign{}
Input length & Original extra bytes & Corrected extra bytes \\
\midrule\noalign{}
\endhead
\bottomrule\noalign{}
\endlastfoot
1,000 & 17,160 & 1,022 \\
10,000 & 160,992 & 832 \\
100,000 & 1,601,000 & 840 \\
\end{longtable}

These observations are consistent with the explicit allocation argument
in Section 7. Their precise byte counts depend on the runtime and do not
establish a machine-independent empirical complexity law. The auxiliary
bound for the corrected algorithm follows from object accounting in
Theorem 1.

\hypertarget{discussion-and-limitations}{%
\section{9 Discussion and
limitations}\label{discussion-and-limitations}}

The omission representation is useful because it changes the number of
reconstruction decisions. A direct feasibility decision for each mutable
index makes \(O(H)\) oracle calls. This reconstruction uses
\(O(\sqrt H\log(h+2))\) calls while retaining scalar states. Its cost
includes an additional logarithm relative to the repository's stated G8
sparse-witness time bound. In that comparison, the G8 analysis states
time \(O(n+H^{5/2}\log(h+2))\) and auxiliary space \(O(H^2\log(h+2))\),
whereas the present theorem removes the logarithm from auxiliary space
at the cost of another logarithm in time (Saridakis 2026, G8 theorem
laboratory). This is a comparison with a repository-stated theoretical
bound, not independent certification of every allocation in the G8
implementation, and not an external state-of-the-art ranking.

Five conclusions must be kept separate. Mathematical correctness is
supported by the explicit derivation. Implementation correctness
includes a verified space defect in the original and a corresponding
correction. The asymptotic contribution is a specialized space-time
tradeoff. Literature priority is unresolved. Scientific significance is
bounded by the narrow regime, established ingredients, and absence of
demonstrated applications or lower bounds. None of these conclusions
follows automatically from any of the others.

Several limitations remain. The algorithm assumes \(h<n\) and does not
solve all modulus regimes within these bounds. Arithmetic on
arbitrary-size values is charged through a wide-word model. The
experiments are finite, predominantly small, and not an industrial
benchmark. The independent reference programs and the proof were
prepared in an AI-assisted workflow within the same session; external
expert replication remains valuable. The literature search is
source-specific but not exhaustive. The correctness theorem can survive
a priority discovery even if the paper's novelty case does not.

Useful subsequent work would identify a substantive application or a
broader parameter regime, prove stronger tradeoffs, or establish a
meaningful lower bound. Such possibilities are research questions, not
results of this note. A venue decision should depend on the surviving
specialized theorem and the reproducible correction, not on the
availability of polished manuscript files.

\hypertarget{conclusion}{%
\section{10 Conclusion}\label{conclusion}}

In the high-modulus small-deficit regime, exact modular feasibility
confines all omissions to a short prefix and bounds their number by a
square root of the deficit. A constrained sorted-difference scalar DP
then supports monotone searches for the lexicographically greatest
omitted tuple, whose complement is the globally least complete optimal
witness. Charging all searches and output operations gives quadratic
auxiliary wide-word space and \(O(n+(h+1)^{5/2}\log^2(h+2))\) time for
the corrected algorithm. The result is a specialized composition of
established methods with a documented implementation correction; its
literature priority and publication significance remain open to further
assessment.

\hypertarget{artifact-availability-and-disclosure}{%
\section{Artifact availability and
disclosure}\label{artifact-availability-and-disclosure}}

The reproducible audit snapshot is linked in the references (Saridakis
2026). From that commit, run
\texttt{python\ research/v03/g9\_independent\_audit/audit.py} and
\texttt{python\ -m\ unittest\ tests.test\_v03\_g9\_quadratic\_lex\ -v}
with assertions enabled. The first command regenerates
\texttt{results.json}; use a disposable checkout to preserve the
committed evidence. The companion manuscript folder includes the
claim-evidence map, source registry, review, contribution stress test,
evidence hashes, and build instructions. Original G3--G9 research files
are preserved.

This manuscript was drafted with AI assistance, including source
inspection, mathematical reconstruction, code execution, literature
searches, and adversarial review. That review is an internal assisted
critique, not external peer review. The named author's review and
responsibility for any eventual submission are required; no affiliation,
funding award, external endorsement, or venue acceptance is asserted.

\hypertarget{references}{%
\section*{References}\label{references}}
\addcontentsline{toc}{section}{References}

\hypertarget{refs}{}
\begin{CSLReferences}{1}{0}
\leavevmode\vadjust pre{\hypertarget{ref-erickson2019}{}}%
Erickson, Jeff. 2019. \emph{Algorithms}. 1st ed. Self-published.
\url{https://jeffe.cs.illinois.edu/teaching/algorithms/}.

\leavevmode\vadjust pre{\hypertarget{ref-grosse2006}{}}%
Große, André, Jörg Rothe, and Gerd Wechsung. 2006. {``A Note on the
Complexity of Computing the Smallest Four-Coloring of Planar Graphs.''}
\url{https://arxiv.org/abs/cs/0106045v2}.

\leavevmode\vadjust pre{\hypertarget{ref-jin2020}{}}%
Jin, Ce, Nikhil Vyas, and Ryan Williams. 2020. {``Fast Low-Space
Algorithms for Subset Sum.''} \url{https://arxiv.org/abs/2011.03819}.

\leavevmode\vadjust pre{\hypertarget{ref-kluev2013}{}}%
Kluev, Evgeny. 2013. {``Answer to {Longest convex subsequence in an
array}.''} Stack Overflow.
\url{https://stackoverflow.com/questions/20664920/longest-convex-subsequence-in-an-array}.

\leavevmode\vadjust pre{\hypertarget{ref-mocoso2026}{}}%
Saridakis, Stamatis-Christos. 2026. {``{MOCOSO} Research Artifacts and
Independent {G9} Audit.''}
\url{https://github.com/SaridakisStamatisChristos/MOCOSO/tree/43353e2b9a659d736b627788ec0c46554a090c75/research/v03}.

\leavevmode\vadjust pre{\hypertarget{ref-turner2011}{}}%
Turner, Lara. 2011. {``Variants of the Shortest Path Problem.''}
\emph{Algorithmic Operations Research} 6: 91--104.
\url{https://journals.lib.unb.ca/index.php/AOR/article/download/18312/20525/24773}.

\end{CSLReferences}

\end{document}
